% ============================================================
% 分片：s02b-fieldtheory  对应原文 PDF 第 10--16 页（印刷页码 10--16）
% 负责章节：§2 FIELD THEORY BACKGROUND（后半，至第 16 页 §3 标题之前）
% 译者：t03-s2b   审校状态：初稿
% 审校：t04-s3a 已完成逐式核对（修正 4 处）
% ============================================================
% 边界：起于第 10 页第一词（承接第 9 页 (2.16) 关于扭手征超场 $\Sigma$ 分量展开的括注）；
% 止于第 16 页 §3 标题之前的两段。与 s03a-model 的约定一致：该片自 §3 标题起，
% 故第 16 页 §3 标题之前的 §2 结尾两段归本分片。
% 公式一律以 pages/page-10.png--page-16.png 为唯一依据逐式核对；文字层仅供措辞参考。
% 原文的非编号小标题（Lagrangians / Twisted Chiral Superpotential / Symmetries）
% 按 SPEC §3 用 \subsection*{} 保持层级、不加编号。
% ==== 原文 PDF 第 10 页（印刷页 10）====
\label{op:10}
（这些涉及导数的复杂项由前三项借助扭手征条件 $\bar{\mathcal D}_+\Sigma=\mathcal D_-\Sigma=0$ 确定；它们可以通过 $x^m\to y^m+i\theta^\alpha\sigma^m_{\alpha\dot\alpha}\bar\theta^{\dot\alpha}$ 的扭化版本消去。）

% [Tutor Note: 推导补足 —— 扭手征超场的定义是 $\bar{\mathcal D}_+\Sigma=\mathcal D_-\Sigma=0$，而手征超场满足的是 $\bar{\mathcal D}_\pm\Phi=0$；这两个条件各自只涉及一半的费米型导数，是二维 $N=2$ 超空间（而非四维）才允许的结构。由该约束，$\Sigma$ 的 $\theta$ 展开中只有 $\theta^+$ 与 $\bar\theta^-$ 是独立的奇坐标，其余分量全部由 $\sigma$ 的导数确定：(2.16) 中带 $\partial_0\pm\partial_1$ 的项就是这样产生的，这正是括注所说“涉及导数的复杂项由前三项确定”的含义。代换 $x^m\to y^m+i\theta^\alpha\sigma^m_{\alpha\dot\alpha}\bar\theta^{\dot\alpha}$ 是把普通坐标换成扭手征坐标，使 $\Sigma$ 只依赖于 $y^m,\theta^+,\bar\theta^-$，因而 (2.26)、(2.29) 的测度只取 $\mathrm{d}\theta^+\mathrm{d}\bar\theta^-$。]

\subsection*{拉格朗日量}

在超空间中写出拉格朗日量是直接的。我们只考虑超可重整（super-renormalizable）理论；它们将是其普适类中特别简单的代表。

这些拉格朗日量具有如下形式
\begin{equation}\tag{2.17}
L=L_{kin}+L_W+L_{gauge}+L_{D,\theta}
\end{equation}
其中四项分别是手征超场的动能、超势相互作用、规范场的动能，以及 Fayet--Iliopoulos 项与 $\theta$ 角；它们都按如下方式构造。

为简单起见，我们针对阿贝尔规范群 $U(1)^s$ 的情形写出公式，它如上所述由矢量超场 $V_a$（$a=1\ldots s$）描述。这将是我们主要关心的情形。设有 $k$ 个电荷为 $Q_{i,a}$ 的手征超场 $\Phi_i$。$\Phi_i$ 可解释为 $Z=\mathbb{C}^k$ 的一个副本上的坐标；其动能由 $Z$ 上的凯勒度规决定。为满足超可重整性，该度规应当是平坦的。在适当选取坐标后，对应于这种度规的拉格朗日量为
\begin{equation}\tag{2.18}
L_{kin}=\int \mathrm{d}^{2}y\,\mathrm{d}^{4}\theta\sum_i\bar\Phi_i\Phi_i
=\int \mathrm{d}^{2}y\,\mathrm{d}^{4}\theta\sum_i\bar\Phi_{0,i}e^{2\sum_a Q_{i,a}V_a}\Phi_{0,i}.
\end{equation}
（$\Phi_{0,i}$ 按 (2.13) 由 $\Phi_i=\exp\bigl(\sum_a Q_{i,a}V_a\bigr)\Phi_{0,i}$ 定义。）写成分量形式即得

% ==== 原文 PDF 第 11 页（印刷页 11）====
\label{op:11}
% [跨页接续: 原文第 10 页末 "In components" 与第 11 页首 "this becomes" 属同一句，译文按语义完整保留。]

\begin{equation}\tag{2.19}
\begin{aligned}
L_{kin}={}&\sum_i\int \mathrm{d}^{2}y\,\Bigl(-D_\rho\bar\phi_i D^\rho\phi_i+i\bar\psi_{-,i}(D_0+D_1)\psi_{-,i}+i\bar\psi_{+,i}(D_0-D_1)\psi_{+,i}+|F_i|^2\\
&-2\sum_a\bar\sigma_a\sigma_a Q_{i,a}^2\bar\phi_i\phi_i-\sqrt{2}\sum_a Q_{i,a}\bigl(\bar\sigma_a\bar\psi_{+,i}\psi_{-,i}+\sigma_a\bar\psi_{-,i}\psi_{+,i}\bigr)+\sum_a D_a Q_{i,a}\bar\phi_i\phi_i\\
&-\sum_a i\sqrt{2}Q_{i,a}\bar\phi_i\bigl(\psi_{-,i}\lambda_{+,a}-\psi_{+,i}\lambda_{-,a}\bigr)-\sum_a i\sqrt{2}Q_{i,a}\phi_i\bigl(\bar\lambda_{-,a}\bar\psi_{+,i}-\bar\lambda_{+,a}\bar\psi_{-,i}\bigr)\Bigr)
\end{aligned}
\end{equation}

% [Tutor Note: 易错预警 —— (2.19) 的全部符号由世界面号差 $(-,+)$ 与 (2.23) 的约定固定：$D_0\pm D_1$ 是光锥方向的协变导数，$\psi_{+,i}$ 只与 $D_0-D_1$ 相配、$\psi_{-,i}$ 只与 $D_0+D_1$ 相配，下标的正负号不可随手互换。核对时注意三点：(i) 标量动能项整体带负号 $-D_\rho\bar\phi_iD^\rho\phi_i$，在 $(-,+)$ 号差下这正是正能条件；(ii) $Q_{i,a}$ 是 $\Phi_i$ 对第 $a$ 个 $U(1)$ 的荷，$\bar\sigma_a\sigma_a Q_{i,a}^2$ 一项来自扭手征超场 $\Sigma_a$ 的标量分量；(iii) 第三行的 $D_aQ_{i,a}\bar\phi_i\phi_i$ 与 (2.23) 的 $\frac{1}{2}D_a^2$ 联立消去 $D_a$，并与 (2.24) 的 $-rD$ 合并，即得 §3 的 (3.1)--(3.2)：$D_a=-e_a^2\bigl(\sum_iQ_{i,a}|\phi_i|^2-r\bigr)$。]

其中世界面度规为 $\mathrm{d}s^2=-(\mathrm{d}y^0)^2+(\mathrm{d}y^1)^2$。

拉格朗日量中涉及手征超场的另一部分，由 $Z$ 上一个规范不变的全纯函数 $W$ 构造，$W$ 称为超势（superpotential）。在我们的模型中，$W$ 总是多项式。拉格朗日量的相应部分为
\begin{equation}\tag{2.20}
L_W=-\int \mathrm{d}^{2}y\,\mathrm{d}\theta^+\mathrm{d}\theta^-\,W(\Phi_i)\big|_{\bar\theta^+=\bar\theta^-=0}-\mathrm{h.c.}
\end{equation}
该积分只对奇数坐标的一半进行；由于 $\Phi_i$ 是手征超场且 $W$ 全纯，这一运算才是超对称的。写成分量形式，
\begin{equation}\tag{2.21}
L_W=-\int \mathrm{d}^{2}y\Bigl(F_i\frac{\partial W}{\partial\phi_i}+\frac{\partial^2 W}{\partial\phi_i\partial\phi_j}\psi_{-,i}\psi_{+,j}\Bigr)-\mathrm{h.c.}
\end{equation}

规范动能由如上定义的扭手征超场 $\Sigma_a$ 构造。引入规范耦合常数 $e_a$（$a=1\ldots s$），规范动能为
\begin{equation}\tag{2.22}
L_{gauge}=-\sum_a\frac{1}{4e_a^2}\int \mathrm{d}^{2}y\,\mathrm{d}^{4}\theta\,\bar\Sigma_a\Sigma_a.
\end{equation}
写成分量形式即
\begin{equation}\tag{2.23}
\begin{aligned}
L_{gauge}={}&\sum_a\frac{1}{e_a^2}\int \mathrm{d}^{2}y\,\Bigl(\frac{1}{2}v_{01,a}^2+\frac{1}{2}D_a^2+i\bar\lambda_{+,a}(\partial_0-\partial_1)\lambda_{+,a}\\
&+i\bar\lambda_{-,a}(\partial_0+\partial_1)\lambda_{-,a}-|\partial_\rho\sigma_a|^2\Bigr).
\end{aligned}
\end{equation}

涉及规范多重态的一些附加项将起关键作用。

% ==== 原文 PDF 第 12 页（印刷页 12）====
\label{op:12}
% [跨页接续: 原文第 11 页末 "will play an essential" 与第 12 页首 "role" 属同一句，译文按语义完整保留。]

为简单起见，以下公式针对单个规范多重态 $V$、即规范群恰为 $U(1)$ 的情形写出。推广只需像 (2.23) 那样对各分量求和即得。

通过考察超对称变换规律可以看出，Fayet--Iliopoulos 项
\begin{equation}\tag{2.24}
-r\int \mathrm{d}^{2}y\,D
\end{equation}
是超对称的。它通常写成超空间中的 $-\int \mathrm{d}^{2}y\,\mathrm{d}^{4}\theta\,V$，但这一表示并不令人满意，因为 $V$ 不是规范不变的。稍后将给出一个更好的版本。另一个重要的耦合是 $\theta$ 角，
\begin{equation}\tag{2.25}
\frac{\theta}{2\pi}\int \mathrm{d}v=\frac{\theta}{2\pi}\int \mathrm{d}^{2}y\,v_{01}.
\end{equation}

% [Tutor Note: 推导补足 —— 为什么 $-\int \mathrm{d}^{2}y\,\mathrm{d}^{4}\theta\,V$ 是“不令人满意的表示”：在 $V\to V+\Lambda+\bar\Lambda$（$\Lambda$ 手征）下，被积函数本身发生改变，只因 $\Lambda$ 的最高分量是全导数才使积分不变；被积函数不是规范不变的，就无法用它乘以规范不变算符来构造更一般的耦合。改用规范不变的扭手征超场 $\Sigma=\frac{1}{2\sqrt2}\{\bar{\mathcal D}_+,\mathcal D_-\}$ 后，FI 项与 $\theta$ 角可以合并为 (2.27)：由全纯参数 $t=ir+\theta/(2\pi)$ 的线性扭超势同时给出这两个耦合，$r$ 进入 $t$ 的虚部、$\theta$ 进入实部。(2.26) 的两个公式正是把 $(D\mp iv_{01})$ 写成扭手征测度下的分量积分。]

$\theta$ 应视为角变量（物理量在 $\theta$ 中以 $2\pi$ 为周期，但可能并不关于 $\theta$ 光滑、甚至不连续），因为量 $\int \mathrm{d}v/2\pi$ 度量阿贝尔规范场的第一陈类，并且总是整数（在适当的边界条件下）。

上述相互作用可以借助扭手征超场 $\Sigma$ 在超空间中简洁地写出。事实上，
\begin{equation}\tag{2.26}
\begin{aligned}
\int \mathrm{d}^{2}y\,\mathrm{d}\theta^+\mathrm{d}\bar\theta^-\,\Sigma\big|_{\theta^-=\bar\theta^+=0}&=\sqrt{2}\int \mathrm{d}^{2}y\,(D-iv_{01})\\[2pt]
\int \mathrm{d}^{2}y\,\mathrm{d}\theta^-\mathrm{d}\bar\theta^+\,\bar\Sigma\big|_{\theta^+=\bar\theta^-=0}&=\sqrt{2}\int \mathrm{d}^{2}y\,(D+iv_{01}).
\end{aligned}
\end{equation}
由 (2.26) 可知，可以写出
\begin{equation}\tag{2.27}
\begin{aligned}
L_{D,\theta}&=\int \mathrm{d}^{2}y\Bigl(-rD+\frac{\theta}{2\pi}v_{01}\Bigr)\\
&=\frac{it}{2\sqrt{2}}\int \mathrm{d}^{2}y\,\mathrm{d}\theta^+\mathrm{d}\bar\theta^-\,\Sigma\big|_{\theta^-=\bar\theta^+=0}-\frac{i\bar t}{2\sqrt{2}}\int \mathrm{d}^{2}y\,\mathrm{d}\theta^-\mathrm{d}\bar\theta^+\,\bar\Sigma\big|_{\theta^+=\bar\theta^-=0}
\end{aligned}
\end{equation}
其中
\begin{equation}\tag{2.28}
t=ir+\frac{\theta}{2\pi}.
\end{equation}

% [Tutor Note: 物理图像 —— (2.28) 的 $t=ir+\theta/(2\pi)$ 是二维规范理论中 FI 参数与 $\theta$ 角的复化组合，即复化凯勒类（complexified Kähler class）：$r$ 是凯勒类的大小，$\theta$ 是与它配对、相当于 $B$ 场的角变量。理论的全部全纯数据都通过 $t$ 出现，这正是镜像对称中复化凯勒模的雏形；后文 §3 把 $t$ 当作 GLSM 唯一的扭超势参数使用。]

% ==== 原文 PDF 第 13 页（印刷页 13）====
\label{op:13}

\subsection*{扭手征超势}

% [背景补充: 四维中超对称规范场的场强是手征超场 $W_\alpha$，而“扭手征”是二维特有的结构——四维没有与 $\Sigma$ 对应的扭手征超场，因此四维理论没有“扭超势”这一结构。只有在二维，规范场强 $\Sigma=\frac{1}{2\sqrt2}\{\bar{\mathcal D}_+,\mathcal D_-\}$ 恰为扭手征超场，才允许写出全纯的扭超势。这就是原文所说“在四维中没有相近的类比”的确切含义。]

(2.27) 有一个推广，它可能并不熟悉，因为它在四维中没有相近的类比。（不过参见 [21,20] 中密切相关的讨论。）可以引入全纯的“扭超势”（twisted superpotential）$\widetilde W(\Sigma)$，并写出
\begin{equation}\tag{2.29}
\Delta L=\int \mathrm{d}^{2}y\,\mathrm{d}\theta^+\mathrm{d}\bar\theta^-\,\widetilde W(\Sigma)\big|_{\theta^-=\bar\theta^+=0}+\mathrm{h.c.}
\end{equation}
写成分量形式，即
\begin{equation}\tag{2.30}
\Delta L=\int \mathrm{d}^{2}y\Bigl(\sqrt{2}\widetilde W'(\sigma)(D-iv_{01})+2\widetilde W''(\sigma)\bar\lambda_+\lambda_-\Bigr)+\mathrm{h.c.}
\end{equation}

% [Tutor Note: 推导补足 —— 由 (2.29) 到 (2.30)：$\int\mathrm{d}\theta^+\mathrm{d}\bar\theta^-$ 取出 $\widetilde W(\Sigma)$ 的 $\theta^+\bar\theta^-$ 分量。用 (2.16) 的展开 $\Sigma=\sigma-i\sqrt2\theta^+\bar\lambda_+-i\sqrt2\bar\theta^-\lambda_-+\sqrt2\theta^+\bar\theta^-(D-iv_{01})+\ldots$ 代入 $\widetilde W$ 并按 $\theta^+\bar\theta^-$ 取系数，即得 $\sqrt2\widetilde W'(\sigma)(D-iv_{01})+2\widetilde W''(\sigma)\bar\lambda_+\lambda_-$。$\Sigma$ 分量中的 $\sqrt2$ 给出前者的 $\sqrt2$，Grassmann 展开的二阶项给出后者的 $2$；若漏掉 $\widetilde W''$ 项，就无法从 (2.30) 得到 (2.31) 中含 $\bar\lambda_+\lambda_-$ 的项。]

当 $\widetilde W$ 取线性函数 $\widetilde W(x)=itx/2\sqrt2$ 时，它恰好约化为 $L_{D,\theta}$。

在本文中，假定微观的 $\widetilde W$ 函数是线性的，部分原因是要保证下面即将讨论的 $R$ 不变性。即便如此，也会发现量子修正可以生成更复杂的扭超势。以下的评述将为这一讨论提供有用的背景。

考虑扭超势 $\widetilde W(x)=-x\ln x\cdot p/\sqrt2$（$p$ 为常数）。计算得
\begin{equation}\tag{2.31}
\Delta L=-p\int \mathrm{d}^{2}y\Bigl((\ln\sigma+1)(D-iv_{01})+\sqrt{2}\,\frac{1}{\sigma}\bar\lambda_+\lambda_-\Bigr)-\mathrm{h.c.}
\end{equation}
由于 $\ln\sigma$ 只在模 $2\pi i\mathbb{Z}$ 的意义下良定义，(2.31) 作为实值泛函并非良定义。然而，在量子力学中，只要 $\Delta L$ 在模 $2\pi$ 意义下良定义就足够了。由于 $\int \mathrm{d}^{2}y\,v_{01}$ 取值于 $2\pi\mathbb{Z}$，$\Delta L$ 在模 $2\pi$ 意义下良定义当且仅当
\begin{equation}\tag{2.32}
4\pi p\in\mathbb{Z}.
\end{equation}

% [Tutor Note: 推导补足与物理图像 —— 把 (2.31) 按 $D$ 与 $v_{01}$ 展开：$\Delta L=-2p(\ln|\sigma|+1)D-2p(\arg\sigma)\,v_{01}+\ldots$，故它相当于把 FI 参数与 $\theta$ 角分别移动 $r\to r+2p(\ln|\sigma|+1)$、$\theta\to\theta-4\pi p\arg\sigma$。当 $\sigma$ 绕原点一周时 $\ln\sigma\to\ln\sigma+2\pi i$，即 $\theta\to\theta-8\pi^2p$；作用量单值性要求该移动为 $2\pi$ 的整数倍，这正是 (2.32) 的 $4\pi p\in\mathbb{Z}$。这也是二维轴反常的体现：手征旋转使 $\theta$ 平移 $2\pi\times$(反常系数)，而 $\ln\sigma$ 的多值性恰好实现了这一平移。]

由于 (2.31) 在 $\sigma=0$ 处具有奇异性，它若由量子修正生成，则只可能出现在有效性被限制于大 $\sigma$ 的近似之中。

% ==== 原文 PDF 第 14 页（印刷页 14）====
\label{op:14}
% [跨页接续: 原文第 13 页末 "an approximation whose validity would" 与第 14 页首 "be limited to large $\sigma$" 属同一句，译文已按语义合并于上段。]

这正是我们将要发现的结果。为将来使用，这里记下：若 $\sigma$ 很大且变化缓慢，则把 (2.31) 加入拉格朗日量将起到把 $r$ 移至一个有效值的作用，
\begin{equation}\tag{2.33}
r_{eff}=r+2p\ln|\sigma|.
\end{equation}

% [Tutor Note: 易错预警 —— (2.33) 只保留 $|\sigma|$ 的依赖：严格展开 (2.31) 的 $D$ 线性项会得到 $r\to r+2p\ln|\sigma|+2p$，其中与场无关的常数 $2p$ 只是 $r$ 的常数重整化，可吸收进 $r$ 的定义，故 (2.33) 写作 $r_{eff}=r+2p\ln|\sigma|$。此处符号是本段最易抄错的地方，它由三处约定共同固定：(2.24) 的 $-rD$、(2.28) 的 $t=ir+\theta/(2\pi)$、以及 (2.31) 中 $-p(\ln\sigma+1)$ 前的负号；任何一处改号都会使 (2.33) 的 $+$ 变成 $-$。另须注意 $\ln\sigma$ 的虚部移动的是 $\theta$ 而不是 $r$（见上一条批注），把 $|\sigma|$ 误写成 $\sigma$ 会同时弄错两者的漂移。]

\subsection*{对称性}

我们所构造的模型具有二维 $N=2$ 超对称，即两个左行超对称（作用在 $\theta^-,\bar\theta^-$ 上）和两个右行超对称（作用在 $\theta^+,\bar\theta^+$ 上）。自然要问：这些模型是否也具有左行与右行的 $R$ 对称性？右行 $R$ 对称性是一个 $U(1)$ 对称性，在其下 $\theta^+\to e^{i\alpha}\theta^+$、$\bar\theta^+\to e^{-i\alpha}\bar\theta^+$，而 $\theta^-,\bar\theta^-$ 不变。左行 $R$ 对称性则服从把 $+$ 与 $-$ 互换的类似条件。

左行与右行 $R$ 对称性之所以重要，是因为它们是 $N=2$ 超共形代数的一部分。若存在，它们将给出关于所讨论模型可能的共形不变极限的重要信息。

从四维的观点看，所讨论的模型中存在这样的左行与右行对称性是相当自然的。许多四维超对称理论会有一个单一的 $R$ 对称性，其中 $\theta^+$ 与 $\theta^-$ 的荷为 $+1$，$\bar\theta^\pm$ 的荷为 $-1$。此外，在约化到二维的过程中，我们可能得到第二个 $U(1)$ 对称性——对应于额外维度的旋转。在此对称性下，$\theta^+$ 与 $\theta^-$ 反向变换。通过取这些操作的适当组合，有望找到左行或右行 $R$ 对称性，在其下 $\theta^+$ 或 $\theta^-$ 应当不变。

在上述二维公式中很容易显式地找到这些对称性。首先略去超势相互作用 $L_W$，并且假定任何存在的扭超势相互作用都是具有 $L_{D,\theta}$ 这一特定形式的线性项。在这些假定下，左行与右行 $R$ 对称性

% ==== 原文 PDF 第 15 页（印刷页 15）====
\label{op:15}
% [跨页接续: 原文第 14 页末 "R symmetries are" 与第 15 页首 "easy to find explicitly" 属同一句，译文按语义完整保留。]

很容易显式写出。例如，存在一个右行 $R$ 对称性（$\epsilon^+=-\bar\epsilon_-$ 的荷为 $1$，$\bar\epsilon^-=\epsilon_+$ 不变），其下非零的荷分配如下：$(\psi_{+,i},F_i,\sigma_a,\lambda_{-,a})$ 的荷为 $(-1,-1,1,1)$；它们的复共轭具有相反的荷；其余场不变。类似地可以构造一个左行 $R$ 对称性（$\epsilon^-=\epsilon_+$ 的荷为 $1$，$\epsilon^+=-\epsilon_-$ 不变），其下 $(\psi_{-,i},F_i,\sigma_a,\lambda_{+,a})$ 的荷为 $(-1,-1,-1,1)$，其余场中性。我们把产生右行与左行对称性的荷分别称为 $J_R$ 与 $J_L$。
% [存疑: 此处原文的四个超对称参数 $(\epsilon^+,\epsilon^-,\epsilon_+,\epsilon_-)$ 在图像上挤在一起，上划线（$\bar\epsilon$ 的横杠）与上标的相对位置难以逐一无歧义判读；译文按图像尽量照录。读者核对时请以原文图像为准。]

最重要的性质之一是 $J_R$ 与 $J_L$ 可能反常。原因在于：例如 $J_R$ 耦合到右行带电费米子 $\psi_{+,i},\bar\psi_{+,i}$，而它所耦合的左行费米子是中性的。使 $J_R$ 的守恒不被规范反常破坏的条件是
\begin{equation}\tag{2.34}
\sum_i Q_{i,a}=0,\qquad \text{for } a=1\ldots s.
\end{equation}

% [Tutor Note: 推导补足 —— (2.34) 是无反常条件。$J_R$ 只与右行费米子耦合，$J_L$ 只与左行费米子耦合，二者与规范场的混合反常系数互为相反数（分别正比于 $\mp\sum_iQ_{i,a}$），故同一条件 (2.34) 同时保证 $J_R$ 与 $J_L$ 守恒。而 $J_L+J_R$ 对左右行费米子一视同仁（$\psi_{+,i}$ 与 $\psi_{-,i}$ 的荷同为 $-1$），是矢量化组合，其规范反常在左右手之间抵消，因此无论 $Q_{i,a}$ 取何值都守恒。]

同样的条件也是使 $J_L$ 确实成为有效对称性所必需的。（无论如何，不论 $Q_{i,a}$ 取何值，和 $J_L+J_R$ 都守恒。）稍后我们将看到，在施加 (2.34) 时还会发生其他各种好的事情。

在继续之前，重要的是指出：所讨论的模型若忽略超势，还具有额外的（无反常的）对称性，它们与左行和右行超对称都对易，因而可以加到这里引入的 $J_L$ 与 $J_R$ 上。它们是手征超场的整体变换，形式为
\begin{equation}\tag{2.35}
\Phi_i\to\exp(i\alpha k_i)\Phi_i,
\end{equation}
其中 $k_i$ 任意；它们与超对称对易，因为假定 $\Phi_i$ 的每个分量都作相同的变换。

% [Tutor Note: 推导补足 —— (2.35) 不改变 $\Phi_i$ 各分量之间的相对权重，故与超对称对易；$k_i$ 任意，给出 $k$ 个额外 $U(1)$（在无超势时）。超势破坏这些对称性：来自 $\int\mathrm{d}^2\theta\,W$ 的耦合在 $J_R$ 下带荷 $1$，要恢复不变性就须让 $W$ 在 (2.35) 下按 $W\to e^{-i\alpha}W$ 变换，即 $W$ 拟齐次。对单项式 $\prod_i\Phi_i^{n_i}$ 而言这一条件就是 $\sum_i n_ik_i=-1$（即加上 $J_R$ 的荷 $1$ 后总荷为零），这正是 §3 中给 $P$ 取 $k_P=-1$、给 $S_i$ 取 $k_{S_i}=0$ 的依据。]

这些 $R$ 对称性在存在超势时是否仍然有效？来自 $\int \mathrm{d}^{2}\theta\,W$ 的耦合在我们迄今为止所定义的 $J_R$ 下具有荷 $1$。

% ==== 原文 PDF 第 16 页（印刷页 16，§2 末段；§3 标题起归 s03a）====
\label{op:16}
% [跨页接续: 原文第 15 页末 "The couplings coming from $\int\mathrm{d}^2\theta\,W$ have charge 1 under $J_R$ as we have defined it" 与第 16 页首 "hitherto." 属同一句；"hitherto" 已并入上句译出，本句止于此处。]

为挽救这一局面，必须给 $J_R$ 加上一个形式为 (2.35) 的变换，在其下 $W\to e^{-i\alpha}W$。若存在这样的 $k_i$ 使 $W$ 按所述方式变换，则称超势 $W$ 是拟齐次的（quasi-homogeneous）。于是，通过给原来的 $J_R$ 加上附加项（并对 $J_L$ 加上类似的附加项），恰在 (2.34) 成立且 $W$ 拟齐次的情形下可以找到左行与右行 $R$ 对称性。

扭超势破坏 $R$ 不变性，除非它取导致我们称之为 $L_{D,\theta}$ 的线性形式。因此，在 $R$ 不变的情形，量子修正不能生成更复杂的扭手征超势。

% [Tutor Note: 物理图像 —— “$R$ 不变 $\Rightarrow$ 不能生成更复杂的扭超势”是 §3 比较 Calabi--Yau 相与 Landau--Ginzburg 相的关键：扭超势只允许全纯项，而微观的线性 $\widetilde W(\Sigma)=\frac{it}{2\sqrt2}\Sigma$ 给出 (2.27)，其 $\theta^+\bar\theta^-$ 分量为 $D-iv_{01}$，在 $J_R$ 下中性；$\Sigma$ 的高次项则会带上 $\sigma$ 的 $R$ 荷，无法保持不变。于是 $R$ 不变模型中的扭超势只有参数 $t$，在比较两种相时可以把 $t$（即复化凯勒模）当作唯一的扭超势数据。]
